\documentclass[12pt,a4paper]{article}
\usepackage[utf8]{inputenc}
\usepackage{graphicx}
\usepackage{hyperref}
\usepackage{geometry}
\usepackage{float}

\geometry{margin=2.5cm}

\hypersetup{
    colorlinks=true,
    urlcolor=blue,
}

\begin{document}

\vspace*{-3cm}
\begin{center}
    {\LARGE \textbf{FIT5032 - Assessed Lab 10}} \\
    \vspace{0.3cm}
    {\large Zi Tao} \\
    {\large Student ID: 36668680}
\end{center}

\section*{Task 10.1: Weather API and CountBook API}

\textbf{Screenshot 10.1a:} WeatherView code in VS Code and browser showing current location weather.

\begin{figure}[H]
    \centering
    \includegraphics[width=\textwidth]{screenshot10_1a.png}
\end{figure}

\newpage

\textbf{Screenshot 10.1b:} CountBookAPI code in VS Code and browser showing JSON data with authors count and total books.

\begin{figure}[H]
    \centering
    \includegraphics[width=\textwidth]{screenshot10_1b.png}
\end{figure}

\newpage

\section*{Task 10.2: Search Weather by City and GetBookAPI}

\textbf{Screenshot 10.2a:} WeatherView code in VS Code and browser showing ``Clayton, AU'' weather search result.

\begin{figure}[H]
    \centering
    \includegraphics[width=\textwidth]{screenshot10_2a.png}
\end{figure}

\newpage

\textbf{Screenshot 10.2b:} Browser showing all books in JSON format on ``GetBookAPI'' page.

\begin{figure}[H]
    \centering
    \includegraphics[width=\textwidth]{screenshot10_2b.png}
\end{figure}

\newpage

\textbf{Screenshot 10.2c:} GitHub commit history showing project updates.

\begin{figure}[H]
    \centering
    \includegraphics[width=\textwidth]{screenshot10_2c.png}
\end{figure}


\section*{GitHub Repository}
Lab 10 Repository: \href{https://github.com/ztao0014/FIT5032_2026_Zi-Tao_36668680_Lab10}{https://github.com/ztao0014/FIT5032\_2026\_Zi-Tao\_36668680\_Lab10}

\section*{Peer Engagement}
I engaged with my peers through in-person discussions during the lab session about using third-party APIs (OpenWeatherMap) with axios and building local API endpoints with Vue.js. I also participated in the Ed Discussion forum to help classmates with API integration and JSON data handling.

\end{document}
